\begin{textsample}{POS Dim 2 – 2009\_c5 – Score -8.00 – 2009\_c5/enade\_2009\_ci\_econ\_d1cc84d7e1.txt}  \label{ex:f2_pos_421}
Capítulo 5 Distribuição dos Conceitos 5.1 Panorama nacional da distribuição dos conceitos A Tabela 5.1 apresenta a quantidade e o percentual de cursos de CIÊNCIAS ECONÔMICAS participantes do ENADE/2009 que se situaram em cada uma das faixas de conceitos, distribuídos de acordo com a Região em que se localizam os cursos das instituições de educação superior, mantenedoras dos cursos responsáveis pela oferta.
O Gráfico 5.1, complementar à Tabela 5.1, apresenta a mesma informação sob a forma de diagrama de caixa para uma melhor visualização da dispersão dos conceitos obtidos.
No diagrama de caixa, a caixa propriamente dita fica limitada pelos quartis.
A barra dentro da caixa em negrito representa a mediana.
Linhas verticais unem os quartis aos valores extremos.
Se estes valores extremos se configurarem como valores atípicos, estes aparecem como pontos isolados e as linhas terminam no maior ( ou menor ) valor ainda não considerado atípico.
Eventuais valores atípicos aparecem como pequenos círculos com a Categoria Administrativa correspondente.
Valores considerados muito atípicos foram representados por estrelas.
Um diagrama de caixa para poucos dados ou com dados muito concentrados pode ter alguns pontos coincidentes.
Por exemplo, se mais de metade dos cursos de uma Região receberem conceito 5, o valor extremo superior, o quartil superior e a mediana coincidirão.
Se o conjunto só contém um dado este apareceria denotado como a barra da mediana.
Observando -se os dados da Tabela 5.1, percebe -se que, em termos nacionais, predominaram os cursos que receberam conceitos 3 e 2, correspondendo respectivamente a 38,4 \% e a 25,9 \% do total nacional.
Exatamente 3,2 \% dos cursos foram classificados com conceito 1.
Conceitos 4 e 5 foram alocados a, respectivamente, 15,7 \% e 5,1 \%.
Dos 216 cursos, 25 ficaram sem conceito, correspondendo a 11,6 \% do total.
Considerando -se separadamente as cinco Regiões ( ver também Gráfico 5.1 ), observou -se que, na Região Norte, 15 cursos participaram dos quais 3 cursos ( 20,0 \% ) obtiveram o conceito mínimo e nenhum obteve o conceito máximo.
Como o conceito 2 foi modal e centralizado, cursos com conceitos abaixo e acima deste valor aparecem como valores atípicos.
Esta Região foi a que apresentou cursos com a pior distribuição de conceitos.
Na Região Nordeste, 38 cursos participaram, dos quais, 6 ( 15,8 \% ) ficaram sem conceito.
Nenhum curso desta Região obteve conceito máximo e 1 curso ( 2,6 \% ) obteve conceito mínimo.
O conceito modal foi 3, obtido por 14 cursos, equivalentes a 36,8 \% do total da Região.
O conceito 2 foi obtido por 12 cursos ( 31,6 \% ) e 5 cursos ( 13,2 \% ) obtiveram o conceito 4.
Na Região Sudeste, 89 cursos participaram, dos quais, a categoria modal foi o conceito 3, obtido por 37 cursos, equivalentes a 41,6 \% do total regional.
O conceito máximo foi alcançado por 10 cursos ( 11,2 \% ), apenas 1 curso ( 1,1 \% ) obteve conceito mínimo, 15 cursos ( 16,9 \% ) obtiveram o conceito 2.
Nesta Região, 12 cursos ( 13,5 \% do total ) ficaram sem conceito.
Esta Região foi a que apresentou cursos com a melhor distribuição de conceitos.
Na Região Sul, 23 dos 51 cursos avaliados, equivalentes a 45,1 \% do total regional, classificaram -se com conceito modal 3.
Nem o conceito máximo nem o conceito mínimo foram alcançados por nenhum curso.
Ficaram sem conceito 4 cursos ( 7,8 \% ), 13 cursos obtiveram conceito 2 ( 25,5 \% ) e 11 cursos obtiveram conceito 4 ( 21,6 \% ).
Na Região Centro-Oeste, 23 cursos participaram, sendo que 30,4 \% ( 7 cursos ) obtiveram conceito modal 2.
O conceito máximo foi alcançado por somente 1 curso ( 4,3 \% ), 2 cursos ( 8,7 \% ) obtiveram conceito mínimo e 3 ( 13,0 \% ) cursos ficaram sem conceito.
O conceito 3 foi atribuído a 6 cursos ( 26,1 \% ) e o conceito 4 por 4 cursos ( 17,4 \% do total ). 5.2 Conceitos por Categoria Administrativa e por Região A Tabela 5.2 apresenta a distribuição dos cursos participantes do ENADE/2009 de CIÊNCIAS ECONÔMICAS, por Categoria Administrativa, de acordo com os conceitos por eles alcançados, segundo Região.
O Gráfico 5.2 apresenta a mesma informação sob a forma de diagrama de caixa para uma melhor visualização da dispersão dos conceitos obtidos.
Somente eventuais valores atípicos aparecem com o identificador da Região na qual o curso foi habilitado.
De acordo com as informações apresentadas, a maioria, 61,1 \%, é de cursos em instituições privadas, 20,8 \% dos cursos eram vinculados a instituições federais, 13,4 \% a instituições estaduais e os restantes 4,6 \% a instituições municipais.
As instituições federais foram as que apresentaram, em termos nacionais, maior proporção de cursos com conceito 5, equivalentes a 8,9 \% do total dessa categoria.
O conceito modal para esta categoria foi 3, com 31,1 \% dos cursos.
Nenhum curso em instituição federal recebeu o conceito mínimo e 25 cursos ( 20,0 \% ) ficaram sem conceito.
Entre as instituições estaduais, 48,3 \% receberam conceito modal 3.
Nenhum curso em instituição estadual recebeu nem o conceito mínimo nem o máximo, sendo que 2 cursos ( 6,9 \% ) ficaram sem conceito.
Entre as instituições municipais, os conceitos modais foram 2 e 3, para 4 cursos cada, ambos com 40,0 \%.
Entre as instituições municipais, 1 ( 10,0 \% ) curso recebeu o conceito mínimo e nenhuma recebeu o conceito máximo.
Os cursos nas instituições privadas tiveram como conceito modal o valor 3 com uma incidência de 38,6 \% seguido pelo conceito 2 ( 29,5 \% ).
Somente 7 cursos ( 5,3 \% ) receberam o conceito máximo, 6 ( 4,5 \% ) receberam o conceito mínimo.
Ficaram sem conceito 14 cursos em instituições privadas ( 10,6 \% ).
Considerando -se a análise por Região observou -se que em todas as Regiões brasileiras, há representação de cursos em todas as categorias administrativas analisadas, a não ser na Região Norte na qual não existem nem instituições estaduais, nem municipais oferecendo cursos na área de CIÊNCIAS ECONÔMICAS.
Instituições privadas constituem a maioria em todas as Regiões, a não ser na Região Nordeste, onde a maioria são as instituições federais.
Na Região Norte, o conceito modal para cursos em instituições federais foram os conceitos 2 e 3 ( ambos com 50,0 \% de incidência ) e nas privadas foi o conceito 2 ( 66,7 \% ).
Como já comentado, nesta Região não há representação de cursos nem em instituição estadual, nem em municipal.
Nenhum curso alcançou o conceito máximo e nem ficou sem conceito.
O conceito mínimo foi alocado somente a instituições privadas ( 33,3 \% ).
Na Região Nordeste, o conceito modal para cursos em instituições municipais e privadas foi o conceito 2 ( respectivamente 100,0 \% e 50,0 \% ).
Nas instituições federais e estaduais o conceito modal foi 3 ( respectivamente 40,0 \% e 75,0 \% ).
Nesta Região, nenhum curso na Região obteve o conceito máximo e apenas 1 curso, de instituição privada ( 7,1 \% ) obteve conceito mínimo.
Ficaram sem conceito 6 cursos, 4 de instituições federais e 2 de instituições privadas.
Na Região Sudeste, o conceito modal para cursos em instituições municipais e privadas foi o conceito 3 ( respectivamente 80,0 \% e 45,6 \% ).
Nas federais o conceito modal foi 4 ( 38,5 \% ) e nas estaduais os conceitos 2, 3 e 4 foram igualmente distribuídos ( 33,3 \% cada ).
Nesta Região, 12 cursos ficaram sem conceito, sendo 4 em instituições federais e 8 em instituições privadas.
Dez cursos alcançaram o conceito máximo : 3 em instituições federais e 7 em privadas.
O mínimo foi alocado somente a um curso em uma instituição privada.
Na Região Sul o conceito modal para cursos em instituições federais ficou dividido entre os conceitos 3 e 4 igualmente ( 50 \% cada ).
Nas instituições estaduais e privadas o conceito modal foi o conceito 3 ( 38,5 \% e 51,7 \% respectivamente ), nas municipais foi o conceito 2 ( 66,7 \% ).
Nesta Região, 4 cursos ficaram sem conceito, 2 de instituições privadas e 2 de instituições estaduais.
Nenhum curso tampouco alcançou nem o conceito máximo nem o conceito mínimo.
Na Região Centro-Oeste, o conceito modal para cursos em instituições federais foi o conceito 4 ( 40,0 \% ).
Nas estaduais, a categoria modal ficou dividida entre os conceitos 2 e 3 com 40 \% cada.
Apenas um curso em instituição municipal foi avaliado, o qual recebeu o conceito mínimo, 1.
Entre as privadas, o conceito modal foi 2 ( 41,7 \% ).
Nesta Região, 3 cursos ficaram sem conceito, 2 de instituições privadas e 1 de instituição federal.
Somente um curso em uma instituição federal alcançou o conceito máximo nesta Região.
O mínimo foi alocado somente a instituições municipais e privadas, 1 curso em cada. 5.3 Conceitos por Organização Acadêmica e por Região Na Tabela 5.3 encontra -se a distribuição dos conceitos atribuídos aos cursos participantes do ENADE/2009 de CIÊNCIAS ECONÔMICAS, por Organização Acadêmica, segundo as Regiões.
O Gráfico 5.3 apresenta a mesma informação sob a forma de diagrama de caixa para uma melhor visualização da dispersão dos conceitos obtidos.
Somente eventuais valores atípicos apareceram com o identificador da Categoria Administrativa.
Os cursos desta área estão alocados em Universidades ( 56,5 \% ), em Faculdades ( 29,6 \% ) e em Centros Universitários ( 13,9 \% ).
Não existe presença de cursos nem em Centros Federais de Educação Tecnológica ( CEFET ), nem em Institutos Federais de Educação, Ciência e Tecnologia ( IFECT ).
Nas Universidades e Centros Universitários o conceito modal foi 3 para 41,8 \% e 46,7 \% dos cursos, respectivamente.
Nas Faculdades, o conceito modal foi 2 para 34,4 \% dos cursos.
Nenhum curso em Centro Universitário recebeu nem o conceito mínimo, nem o conceito máximo.
Já nos outros 2 tipos de Organização Acadêmica, alguns cursos obtiveram tanto o conceito mínimo quanto o máximo : em Universidades, respectivamente 1,6 \% e 4,9 \%, e em Faculdades, valores bem maiores, respectivamente 7,8 \% e 7,8 \%.
Dos 216 cursos participantes, 25 ficaram sem conceito : 14 cursos em universidades ( correspondendo a 11,5 \% dos cursos nas universidades ) e 11 cursos em faculdades ( 17,2 \% em faculdades ).
Acadêmica, alguns cursos obtiveram tanto o conceito mínimo quanto o máximo : em Universidades, respectivamente 1,6 \% e 4,9 \%, e em Faculdades, valores bem maiores, respectivamente 7,8 \% e 7,8 \%.
Dos 216 cursos participantes, 25 ficaram sem conceito : 14 cursos em universidades ( correspondendo a 11,5 \% dos cursos nas universidades ) e 11 cursos em faculdades ( 17,2 \% em faculdades ).
Considerando -se separadamente as cinco Regiões, observa -se que, na Região Norte, a categoria modal para Universidades ( 57,1 \% ) e Centros Universitários ( 100,0 \% ) foi o conceito 2.
Já nas Faculdades, o conceito modal foi 1 ( 60,0 \% ), o único tipo de Organização Acadêmica a receber este conceito.
Nenhum curso nesta Região obteve conceito máximo.
Nenhum curso ficou sem conceito.
Na Região Nordeste para 51,9 \% dos cursos em Universidades, o conceito modal foi 3, enquanto para 63,6 \% dos cursos em Faculdades, o conceito modal foi 2.
Ficaram sem conceito 18,2 \% dos cursos de Faculdades e 14,8 \% dos cursos de universidades.
Nenhum curso nesta Região obteve conceito máximo e um curso de Faculdade, correspondente a 9,1 \% desta categoria, obteve o conceito mínimo.
Na Região Sudeste, para cursos na área de CIÊNCIAS ECONÔMICAS em Universidades, Faculdades e Centros Universitários, a categoria modal foi o conceito 3, para 31,7 \%, 44,8 \% e 57,9 \%, respectivamente, dos cursos nestas categorias acadêmicas.
O conceito 2, foi alcançado por 31,6 \% dos cursos de Centros Universitários, 13,8 \% de Faculdades e 12,2 \% de Universidades.
O conceito máximo foi alcançado em Universidades por 12,2 \% e em Faculdades por 17,2 \% dos cursos de cada categoria institucional.
O conceito mínimo foi alocado apenas a 1 curso de Universidade ( 2,4 \% ).
Ficaram sem conceito 12 cursos, 7 de universidades e 5 de faculdades, correspondendo a, respectivamente, 17,1 \% e 17,2 \% dos cursos em cada modalidade de instituição.
Na Região Sul o conceito modal para as Universidades e Centros Universitários foi 3, alocado, para 50,0 \% dos cursos em cada tipo de Organização Acadêmica.
Nas Faculdades o conceito modal foi 2, atribuído a 54,5 \% dos cursos.
O conceito máximo e o conceito mínimo não foram alcançados por nenhum curso.
Ficaram sem conceito, nesta Região, um total de 4 cursos : 2 de universidades e 2 de faculdades, correspondendo, respectivamente, a 5,9 \%, e 18,2 \% dos cursos em cada tipo de instituição.
Na Região Centro-Oeste, o conceito modal para as Universidades foi 3, para 30,8 \% dos cursos.
Nas Faculdades foi o conceito 2, com uma incidência de 37,5 \% e nos Centros Universitários o conceito modal ficou dividido igualmente entre 2 e 4, com 50,0 \% cada.
O conceito máximo foi alcançado por um curso em Universidade ( 7,7 \% ).
O conceito mínimo foi alocado somente a 1 curso em Universidade ( 7,7 \% ) e 1 em Faculdade ( 12,5 \% ).
Ficaram sem conceito, nesta Região, um total de 3 cursos : 1 de universidades e 2 de faculdades, correspondendo a, respectivamente, 7,7 \% e 25,0 \% dos cursos em cada tipo de instituição acadêmica.

% matched lemmas: 
\end{textsample}
